% Draft status: V0
% Evidence status: checked where possible; TODOs mark missing evidence.
\section{Results and Analysis}

\input{revision_main_results_table}

\begin{figure}[htbp]
    \centering
    \includegraphics[width=0.84\linewidth]{revision_accuracy_disparity_tradeoff.png}
    \caption{Accuracy-disparity trade-off across model states. Each point is a dataset-model-state evaluation. The plot shows why high average accuracy and low geographic disparity should be evaluated jointly.}
    \label{fig:revision-accuracy-disparity-tradeoff}
\end{figure}

\noindent\textbf{Overall accuracy across representation paradigms.}
Table~\ref{tab:revision-main-results} shows that fine-tuning improves object accuracy for all 18 dataset-model pairs. On Dollar Street, the largest accuracy increase is for ConvNeXtV2, from 1.6\% pretrained accuracy to 76.7\% fine-tuned accuracy. On GeoDE, the largest increase is for DeiT3, from 1.8\% to 96.3\%. The strongest fine-tuned accuracy overall is ConvNeXtV2 on GeoDE at 96.6\%, and ConvNeXtV2 is also the highest-accuracy fine-tuned model on Dollar Street at 76.7\%. At the paradigm level, the Transformer Vision Paradigm has the highest mean fine-tuned accuracy across the two datasets at 85.4\%, closely followed by the Vision-Language Foundation Paradigm at 85.1\%. This suggests that attention-based and image-language representations are highly competitive in average recognition, but the differences are small enough that they should not be overinterpreted without additional uncertainty analysis.

\begin{figure}[htbp]
    \centering
    \includegraphics[width=0.88\linewidth]{revision_slope_accuracy.png}
    \caption{Pretrained-to-fine-tuned change in object accuracy. All evaluated model-dataset pairs move upward after fine-tuning, but the magnitude differs by dataset and model.}
    \label{fig:revision-slope-accuracy}
\end{figure}

\input{revision_paradigm_summary_table}

\noindent\textbf{Regional disparity beyond average accuracy.}
The central result is that high average accuracy does not guarantee low geographic disparity. ConvNeXtV2 has the highest fine-tuned accuracy on both datasets, but it does not have the lowest disparity. MAE has the lowest fine-tuned disparity on both datasets: 6.9\% on Dollar Street and 1.9\% on GeoDE. By contrast, ConvNeXtV2 has disparity of 8.5\% on Dollar Street and 2.3\% on GeoDE. This pattern suggests that accuracy and regional balance are distinct empirical properties. It is consistent with worst-group evaluation arguments that aggregate performance can obscure weaker-group behavior \citep{sagawa2020groupdro}, but the result here is specific to geographic regions and should be interpreted descriptively.

\begin{figure}[htbp]
    \centering
    \includegraphics[width=0.88\linewidth]{revision_slope_disparity.png}
    \caption{Pretrained-to-fine-tuned change in geographic disparity. All evaluated model-dataset pairs show a larger best-minus-worst regional gap after fine-tuning.}
    \label{fig:revision-slope-disparity}
\end{figure}

The dataset contrast is also clear. Fine-tuned Dollar Street runs have larger geographic disparities than fine-tuned GeoDE runs. The mean fine-tuned disparity is 8.9\% on Dollar Street and 2.3\% on GeoDE. This may indicate that the processed Dollar Street subset exposes stronger region-conditioned variation for the evaluated models. A safer interpretation is that the two datasets measure different forms of geographic visual variation, and that regional robustness should not be inferred from only one dataset.

\input{revision_per_region_finetuned_table}

\begin{figure}[htbp]
    \centering
    \includegraphics[width=0.92\linewidth]{revision_per_region_accuracy_dollar_street.png}
    \caption{Dollar Street fine-tuned per-region accuracy. Europe is the highest-accuracy region and Africa is the lowest-accuracy region for all nine models in the available result files.}
    \label{fig:revision-per-region-dollar-street}
\end{figure}

\begin{figure}[htbp]
    \centering
    \includegraphics[width=0.92\linewidth]{revision_per_region_accuracy_geode.png}
    \caption{GeoDE fine-tuned per-region accuracy. Most models have WestAsia as the highest-accuracy region and SouthEastAsia as the lowest-accuracy region, although ResNet50 has Europe as its best region.}
    \label{fig:revision-per-region-geode}
\end{figure}

\begin{figure}[htbp]
    \centering
    \includegraphics[width=0.48\linewidth]{revision_heatmap_region_accuracy_dollar_street.png}
    \includegraphics[width=0.48\linewidth]{revision_heatmap_region_accuracy_geode.png}
    \caption{Model-by-region fine-tuned accuracy heatmaps for Dollar Street and GeoDE. The heatmaps make the regional ordering visible without reducing the analysis to a single best-worst pair.}
    \label{fig:revision-region-heatmaps}
\end{figure}

\begin{figure}[htbp]
    \centering
    \includegraphics[width=0.48\linewidth]{revision_heatmap_region_drop_dollar_street.png}
    \includegraphics[width=0.48\linewidth]{revision_heatmap_region_drop_geode.png}
    \caption{Per-model regional drops relative to each model's best region. This view emphasizes where each model loses accuracy relative to its own strongest region.}
    \label{fig:revision-region-drop-heatmaps}
\end{figure}

\noindent\textbf{Pretrained versus fine-tuned behavior.}
The fine-tuning comparison separates four possible outcomes. The best outcome would be accuracy improvement with disparity reduction. The second outcome is accuracy improvement with disparity increase. The third is accuracy decrease with disparity reduction. The worst outcome is accuracy decrease with disparity increase. In the current results, every evaluated pair falls into the second category: fine-tuning improves average accuracy and increases measured geographic disparity. This does not prove that fine-tuning causes regional disparity. One possible explanation is that pretrained evaluations are near the floor because of classifier-head and label-mapping constraints, so regional differences are compressed when almost all predictions are wrong. After fine-tuning, the models become useful recognizers, and regional differences become measurable.

The magnitude of disparity increase is larger on Dollar Street than on GeoDE. For example, ResNet50 on Dollar Street increases from 2.3\% pretrained accuracy to 58.9\% fine-tuned accuracy, while disparity increases from 0.6\% to 12.2\%. MAE on GeoDE increases from 3.3\% to 92.9\% accuracy, while disparity changes only from 1.8\% to 1.9\%. These examples illustrate why the direction of fine-tuning change should be read jointly with the absolute accuracy level.

\begin{figure}[htbp]
    \centering
    \includegraphics[width=0.78\linewidth]{revision_paradigm_disparity_boxplot.png}
    \caption{Fine-tuned geographic disparity by representation paradigm. The self-supervised group has the lowest mean disparity in these runs, but the evidence is descriptive and does not establish that self-supervision causes regional robustness.}
    \label{fig:revision-paradigm-disparity}
\end{figure}

\noindent\textbf{Paradigm-level interpretation.}
The paradigm-level summary in Table~\ref{tab:revision-paradigm-summary} suggests several patterns. The Transformer Vision Paradigm has the highest mean fine-tuned accuracy, while the Self-supervised Foundation Paradigm has the lowest mean fine-tuned disparity. The Vision-Language Foundation Paradigm is competitive in accuracy, but CLIP has the largest fine-tuned GeoDE disparity among all models. The Convolutional Vision Paradigm contains both the strongest average performer, ConvNeXtV2, and the weakest fine-tuned performer, ResNet50, showing that broad family labels can hide important within-family differences.

These patterns may reflect architecture, pretraining objective, data exposure, optimization, or dataset-specific interactions. They should not be read causally. The safer conclusion is that representation paradigms differ in their observed accuracy-disparity profiles, and that regional robustness requires direct measurement rather than inference from model class or average accuracy alone. TODO: add per-class regional gaps and confidence intervals to test whether the same conclusions hold under uncertainty.
